Every \(h\in H_A\) has the form \((d/2,-d/2)\).  Put \(J=\sigma_1^{-2}+\sigma_2^{-2}=r_1/v_1+r_2/v_2\).  The likelihood ratio between positive and negative values of \(d\) is monotone in
\[
 S=X_1/\sigma_1^2-X_2/\sigma_2^2,
\]
and, under gap \(d\), \(2S/J\sim N(d,4/J)\).  The sign rule therefore has regret
\[
 |d|\Phi(-|d|\sqrt J/2).
\]
Maximizing after the substitution \(x=|d|\sqrt J/2\) gives \(2\kappa/\sqrt J\).  For the symmetric two-point prior at gaps \(d=\pm2x_\star/\sqrt J\), where \(x_\star\) maximizes \(x\Phi(-x)\), the likelihood-ratio rule is the sign of \(S\) and its Bayes regret is the same value.  No rule can have smaller maximum risk, proving the exact all-rule formula.

The raw difference \(X_1-X_2\) is \(N(d,\sigma_1^2+\sigma_2^2)\), so the larger-coordinate rule has worst risk \(\kappa\sqrt{\sigma_1^2+\sigma_2^2}\).  Cauchy--Schwarz gives
\[
 (\sigma_1^2+\sigma_2^2)(\sigma_1^{-2}+\sigma_2^{-2})\geq4,
\]
with equality exactly when \(\sigma_1^2=\sigma_2^2\).  This proves the minimax comparison.

For the larger-coordinate benchmark under exact counts, minimizing \(v_1/\alpha_1+v_2/(1-\alpha_1)\) gives the unique stationary point \(\sqrt{v_1}/(\sqrt{v_1}+\sqrt{v_2})\), and the second derivative is positive.  Under Bernoulli labels write \(a=B_{11}\), \(b=B_{12}\), \(c=B_{21}\), \(d_0=B_{22}\), \(A_0=a-b>0\), and \(D_0=d_0-c>0\).  Then
\[
 F_v'(t)=-\frac{v_1A_0}{(b+A_0t)^2}+\frac{v_2D_0}{(d_0-D_0t)^2},
 \qquad
 F_v''(t)=\frac{2v_1A_0^2}{(b+A_0t)^3}+\frac{2v_2D_0^2}{(d_0-D_0t)^3}>0.
\]
The signs at zero and one give the two endpoint conditions, and solving \(F_v'(t)=0\) gives the displayed interior formula.

For the four-node menu, direct binomial evaluation yields the stated matrix, with \(A_0=11/64\) and \(D_0=21/128\).  Substitution in the preceding exact formula gives
\[
 t^\star=-80+\frac{288\sqrt{462}}{77}=0.39383577843\ldots .
\]
Evaluating \(q_1=1/4+(11/64)t^\star\) and \(q_2=3/8-(21/128)t^\star\) gives the stated larger-coordinate regret ratio \(\{(q_1^{-1}+q_2^{-1})/4\}^{1/2}\) and its square as the cluster-count ratio.  For the all-rule game at equal \(v\), minimizing risk is equivalent to maximizing \(q_1+q_2\).  The two column sums are \(81/128\) and \(80/128\), so \(t=1\).  Exact-count rates sum to one, hence the risk ratio is \((81/128)^{-1/2}=\sqrt{128/81}\), and squaring gives \(128/81\).  The efficient-estimator limits and the two LAN lower bounds transfer these Gaussian calculations to the respective local cluster experiments.